\section{形式承诺的四个层级}

在引入符号之前，本章应当先说明每个符号承担的是哪一种承诺。教材中的每个方程，并不都处在同一个地位上。

\begin{enumerate}[nosep]
  \item \textbf{描述性记号。} 用来命名那些在分析上有用的特征，例如即时回报拉力或环境强化。
  \item \textbf{启发式支架。} 为了记账、排序或说明而引入的简单汇总表达。
  \item \textbf{定性单调假设。} 像“更多含糊通常会抬高转变难度”这样的说法，它们看起来合理、在结构上也有用，但仍然不是已校准的定律。
  \item \textbf{一旦夸大就缺乏支持的内容。} 任何关于已识别的普遍系数、唯一正确的函数形式，或本章日常行动模型已有直接实验室测量的说法。
\end{enumerate}

\begin{remark}
对后续章节而言，实践规则很简单：它们可以继承这套记号、这些支架以及这些单调假设，但不能悄悄把这些元素升级成已定论的定量规律。
\end{remark}